\exercisesection{§ 2 的习题}

\begin{enumerate}
\def\labelenumi{\arabic{enumi})}
\item
  设 G 是一个群，B 和 N 是 G 的两个子群，S 是 \(W = N/(B \cap N)\) 的一个子集。对所有 \(w \in W\)，置 \(C(w) = BwB\)。假设第 1 项定义 1 的条件 (T1) 和 (T2) 得到满足，并且对任意 \(s \in S\) 和 \(w \in W\)，两个关系 \(C(s).C(w) = C(sw)\) 与 \(C(s).C(sw) = C(w)\) 至少有一个成立，且 \(B \cup C(s)\) 对所有 \(s \in S\) 都是 G 的子群。证明条件 (T3) 得到满足。此外，若 B 在 \(B \cup C(s)\) 中的指数为 \(\geqslant 3\)，则 (G, B, N, S) 是 Tits 系统。
\item
  设 G 是一个群，B 和 N 是 G 的两个子群，S 是 \(W = N/(B \cap N)\) 的一个子集。令 Z 为包含于 B 中的 G 的正规子群。令 \(B'\) 和 \(N'\) 为 B、N 在 \(G' = G/Z\) 中的自然像。证明从 N 到 \(N'\) 的自然映射定义了从 W 到 \(W' = N'/(B' \cap N')\) 的同构。令 \(S'\) 为 S 在此同构下的像。证明 \((G', B', N', S')\) 是 Tits 系统，当且仅当 \((G, B, N, S)\) 是 Tits 系统。
\item
  设 G 是一个群，B 是 G 的一个子群，\((\mathrm{C}(w))_{w\in\mathbb{W}}\) 是 G 关于 B 的双陪集族。若存在 W 的一个子集 S，使以下条件得到满足，则称 B 为 G 的 Tits 子群：
\end{enumerate}

\begin{enumerate}
\def\labelenumi{(\arabic{enumi})}
\item
  \(C(s)\)（其中 \(s \in S\)）的并集生成 \(G\)；
\item
  对所有 \(s \in S\)，集合 \(B \cup C(s)\) 是 \(G\) 的子群，且 \(B\) 在 \(B \cup C(s)\) 中的指数为 \(\geqslant 3\)；
\item
  对所有 \(s \in S\) 和所有 \(w \in W\)，存在元素 \(w' \in W\)，使得 \(C(s).C(w) \subset C(w) \cup C(w')\)。
\end{enumerate}

从现在起，假设 B 是 G 的 Tits 子群，并给定 W 的一个满足条件 (1)、(2) 和 (3) 的子集 S。

\begin{enumerate}
\def\labelenumi{\alph{enumi})}
\item
  证明 \(C(s)^{-1} = C(s)\) 以及 \(C(s).C(s) = B \cup C(s)\) 对所有 \(s \in S\) 成立。证明对所有 \(s \in S\) 和所有 \(w \in W\)，存在元素 \(w'' \in W\)，使得 \(C(w).C(s) \subset C(w) \cup C(w'')\)。
\item
  若 \(w \in \mathbf{W}\)，则 \(w\) 的长度（记为 \(l(w)\)）是最小的整数 \(n \geqslant 0\)，使得存在 \(s_1, \ldots, s_n \in \mathbf{S}\) 满足 \(\mathrm{C}(w) \subset \mathrm{C}(s_1) \ldots \mathrm{C}(s_n)\)。证明 \(l(w)\) 对所有 \(w \in \mathbf{W}\) 都是有限的。
\end{enumerate}

  设 \(u, v \in W\) 满足 \(l(u) < l(v)\)，并令 \(s \in S\)。证明：若 \(C(v) \subset C(u).C(s)\)（分别，\(C(v) \subset C(s).C(u)\)），则 \(C(v) = C(u).C(s)\)（分别，\(C(v) = C(s).C(u)\)）（按 \(u\) 的长度作归纳。若 \(C(v) \neq C(u).C(s)\)，则 \(C(u).C(s) = C(u) \cup C(v)\)。利用归纳假设，证明存在 \(t \in S\) 和 \(w \in W\)，使得

\[
\mathrm{C} (u) = \mathrm{C} (t). \mathrm{C} (w) \quad \text { with } \quad l (w) = l (u) - 1.
\]

由关系 \(\mathrm{C}(v) \subset \mathrm{C}(u).\mathrm{C}(s) = \mathrm{C}(t).\mathrm{C}(w).\mathrm{C}(s)\) 以及 \(\mathrm{C}(t).\mathrm{C}(w) = \mathrm{C}(u) \neq \mathrm{C}(v)\)，推出存在元素 \(w' \neq w\)，使得

\(C(w') \subset C(w).C(s)\) 且 \(C(v) \subset C(t).C(w')\)，从而 \(l(w') \geqslant l(v) - 1 > l(u) - 1 = l(w)\)。归纳假设现在推出

\[
\mathrm{C} (w ^ {\prime}) = \mathrm{C} (w). \mathrm{C} (s).
\]

此外，

\[
\begin{array}{c} \mathrm{C} (t). \mathrm{C} (u). \mathrm{C} (s) = \mathrm{C} (t). \mathrm{C} (t). \mathrm{C} (w). \mathrm{C} (s) = \mathrm{C} (t). \mathrm{C} (w). \mathrm{C} (s) \cup \mathrm{C} (w). \mathrm{C} (s) \\ = \mathrm{C} (u). \mathrm{C} (s) \cup \mathrm{C} (w ^ {\prime}) = \mathrm{C} (u) \cup \mathrm{C} (v) \cup \mathrm{C} (w ^ {\prime}) \end{array}
\]

又有

\[
\mathrm{C} (t). \mathrm{C} (u). \mathrm{C} (s) = \mathrm{C} (t). (\mathrm{C} (u) \cup \mathrm{C} (v)) \supset \mathrm{C} (t). \mathrm{C} (t). \mathrm{C} (w) \supset \mathrm{C} (w),
\]

由于 \(w \neq u, v, w'\)，这就矛盾了。

\begin{enumerate}
\def\labelenumi{\alph{enumi})}
\setcounter{enumi}{2}
\tightlist
\item
  证明：对所有 \(w \in \mathbf{W}\) 和所有 \(s \in S\)，存在唯一元素，记为 \(s.w\)（分别为 \(w * s\)），异于 \(w\)，并满足
\end{enumerate}

\[
\begin{array}{c} \mathrm{C} (s. w) \subset \mathrm{C} (s). \mathrm{C} (w) \subset \mathrm{C} (w) \cup \mathrm{C} (s. w) \\ (\text {resp.} \mathrm{C} (w * s) \subset \mathrm{C} (w). \mathrm{C} (s) \subset \mathrm{C} (w) \cup \mathrm{C} (w * s)). \end{array}
\]

（按 \(l(w)\) 归纳证明 \(C(s).C(w) \neq C(w)\)。为此，写成

\[
\mathrm{C} (w) = \mathrm{C} (u). \mathrm{C} (t)
\]

其中 \(t \in S\)、\(u \in W\)，且 \(l(w) = l(u) + 1\)。若 \(C(s).C(w) = C(w)\)，则

\[
\mathrm{C} (s). \mathrm{C} (u). \mathrm{C} (t) = \mathrm{C} (u). \mathrm{C} (t)
\]

在右边乘以 \(\mathrm{C}(t)\)，得到 \(\mathrm{C}(u) \cup \mathrm{C}(w) = \mathrm{C}(s).\mathrm{C}(u) \cup \mathrm{C}(w)\)。由于归纳假设给出 \(\mathrm{C}(u) \neq \mathrm{C}(s).\mathrm{C}(u)\)，由 b) 得

\[
\mathrm{C} (s). \mathrm{C} (u) = \mathrm{C} (w),
\]

从而

\[
\mathrm{C} (u) \subset \mathrm{C} (s). \mathrm{C} (s). \mathrm{C} (u) = \mathrm{C} (s). \mathrm{C} (w) = \mathrm{C} (w),
\]

这是荒谬的）。

\begin{enumerate}
\def\labelenumi{\alph{enumi})}
\setcounter{enumi}{3}
\item
  设 \(s \in S\)。证明映射 \(p_s: w \mapsto s.w\)（分别为 \(q_s: w \mapsto w * s\)）是 \(W\) 的一个置换，并且 \(p_s^2 = \operatorname{Id}(\operatorname{resp.} q_s^2 = \operatorname{Id})\)。证明 \(p_s \circ q_t = q_t \circ p_s\) 对所有 \(s, t \in S\) 成立（研究 \(C(s).C(w).C(t)\)（其中 \(w \in W\)）的乘积，并证明 \((s.w) * t \in \{w, s.w, w * t, s.(w * t)\}\)；证明 \((s.w) * t \notin \{s.w, w * t\}\)，并证明若 \((s.w) * t = w\)，则 \(s.w = w * t\) 且 \(w = s.(w * t)\)）。
\item
  证明由 \(p_{s}\)（分别为 \(q_{s}\)）在 \(s \in S\) 中生成的置换群 P（分别为 Q）在 W 上单纯传递地作用（为证明 P 是传递的，使用 (2) 并按第 1 项引理 1 的方法论证；为证明 P 是单纯传递的，使用 d)）。推出 W 有唯一的群结构，使得映射 \(p \mapsto p(e)\)
\end{enumerate}

 （其中 \(e\) 表示 \(W\) 中满足 \(B = C(e)\) 的元素）是从 \(P\) 到 \(W\) 的同构。映射 \(q \mapsto q(e)\) 于是也是同构；此外，\(s.w = w * s = sw\) 对所有 \(s \in S\) 和所有 \(w \in W\) 成立，且 \(C(w)^{-1} = C(w^{-1})\)。

\begin{enumerate}
\def\labelenumi{\alph{enumi})}
\setcounter{enumi}{5}
\item
  证明对 \((W,S)\) 是一个 Coxeter 系统，并推广第 4 项的结果。
\item
  设 X 是 S 的一个子集，\(W_{X}\) 是由 X 生成的 W 的子群。证明并集 \(G_{X}\)，它由 \(C(w)\)（其中 \(w \in W_{X}\)）组成，是 G 的子群，并且第 5 项定理 3 仍然成立。证明 B 是 \(G_{X}\) 的 Tits 子群。推广第 5 项命题 2 和 3，以及第 6 项定义 2、命题 4 和定理 4。证明 S 正是 \(w \in W\) 中满足 \(B \cup C(w)\) 是异于 B 的 G 的子群的那些元素的集合。因此，Coxeter 系统 (W, S) 和群 W（称为 (G, B) 的 Weyl 群）只依赖于 (G, B)。
\item
  设 N 是 G 的一个子群，满足 \(B \cap N\) 在 N 中正规，并且每个关于 B 的双陪集 \(C(w)\) 与 N 的交是关于 \(B \cap N\) 的一个双陪集。证明群 \(N/(B \cap N)\) 可以认作 W，并且 \((G, B, N, S)\) 是一个 Tits 系统。
\end{enumerate}

\begin{enumerate}
\def\labelenumi{\arabic{enumi})}
\setcounter{enumi}{3}
\tightlist
\item
  设 \((G, B, N, S)\) 和 \((G', B', N', S')\) 是两个 Tits 系统，满足 \(G = G'\) 和 \(B = B'\)，其 Weyl 群分别为 W 和 \(W'\)。令 j 为从 W 到 \(W'\) 的双射，由关系
\end{enumerate}

\[
\mathrm{B} w \mathrm{B} = \mathrm{B} ^ {\prime} j (w) \mathrm{B} ^ {\prime}.
\]

证明 \(j\) 是 \(W\) 与 \(W'\) 之间的群同构，且 \(j(S) = S'\)。

\begin{enumerate}
\def\labelenumi{\arabic{enumi})}
\setcounter{enumi}{4}
\tightlist
\item
  设 \(\Sigma = (G, B, N, S)\) 为一个 Tits 系统。置 \(T = B \cap N\)，并令 \(\tilde{N}\) 为 \(N\) 的正规化子。
\end{enumerate}

\begin{enumerate}
\def\labelenumi{\alph{enumi})}
\tightlist
\item
  设 \(b \in B \cap \hat{N}\)。证明 \(bnb^{-1}n^{-1} \in B \cap N\) 对所有 \(n \in N\) 成立（置 \(bn = n'b\)，并使用定理 1），且 b 属于共轭子群的交 \(\tilde{T}\)，其中共轭子群为 \(nBn^{-1}\)（对所有 \(n \in N\)）。证明 \(\tilde{T} \cap N = T\)。
\end{enumerate}

若 \(\tilde{\mathrm{T}} = \mathrm{T}\)，则称 \(\Sigma\) 是饱和的。

\begin{enumerate}
\def\labelenumi{\alph{enumi})}
\setcounter{enumi}{1}
\item
  置 \(\tilde{N} = N.\tilde{T}\)。证明 \(\tilde{N}\) 是 G 的一个子群，包含 \(\tilde{T}\) 这个正规子群，且 \(\tilde{N} \cap B = \tilde{T}\)。证明 N 到 \(\tilde{N}\) 的嵌入定义了从 \(\Sigma\) 的 Weyl 群 W 到 \(\tilde{N}/\tilde{T}\) 的同构。
\item
  证明 \((\mathbf{G},\mathbf{B},\tilde{\mathbf{N}},j(\mathbf{S}))\) 是一个饱和的 Tits 系统，称为与 \(\Sigma\) 相关的 Tits 系统。
\end{enumerate}

\begin{enumerate}
\def\labelenumi{\arabic{enumi})}
\setcounter{enumi}{5}
\item
  沿用第 2 项的记号，并令 \(\mathbf{N}_0\) 为 \(\mathbf{N}\) 的由所有元素均为 0 或 1 的矩阵组成的子群。证明 \(\mathbf{B} \cap \mathbf{N}_0 = \mathbf{T} \cap \mathbf{N}_0 = \{1\}\)，并且自然映射 \(j\) 从 \(\mathbf{N}_0\) 到 \(\mathbf{W} = \mathbf{N} / \mathbf{T}\) 是同构。置 \(S_0 = j^{-1}(S)\)。证明 \((\mathbf{G}, \mathbf{B}, \mathbf{N}_0, S_0)\) 是 Tits 系统，且 \((\mathbf{G}, \mathbf{B}, \mathbf{N}, S)\) 是相关的饱和 Tits 系统。
\item
  设 G 是作用于集合 E 的一个群。若 \(x, y, x', y' \in E\) 满足 \(x \neq y\) 和 \(x' \neq y'\)，都存在元素 \(g \in G\) 使得 \(g.x = x'\) 且 \(g.y = y'\)，则称 G 在 E 上二重传递地作用。
\end{enumerate}

\begin{enumerate}
\def\labelenumi{\alph{enumi})}
\item
  设 \((G, B, N, S)\) 是 Weyl 群 W 阶为 2 的 Tits 系统。证明 G 在 G/B 上二重传递地作用。
\item
  设 G 是二重传递地作用于集合 E 的一个群。假设 \(\operatorname{Card}(E) \geqslant 3\)。令 \(e \in E\)，令 B 为 e 的稳定子群。令 \(x \in E\) 且 \(x \neq e\)，并令 \(n \in G\) 满足 \(n(e) = x\) 和 \(n(x) = e\)。令 N 为由 n 生成的 G 的子群。令 s 为 n 在 N/T 中的自然像。证明 \((G, B, N, \{s\})\) 是 Weyl 群阶为 2 的 Tits 系统。
\end{enumerate}

\begin{enumerate}
\def\labelenumi{\arabic{enumi})}
\setcounter{enumi}{7}
\tightlist
\item
  设 \((\mathrm{G},\mathrm{B},\mathrm{N},\mathrm{S})\) 为一个 Tits 系统；置 \(T=B\cap N\)，W=N/T。令 \(\tilde{G}\) 为包含 G 为正规子群的一个群。假设对所有 \(h\in\tilde{G}\)，都存在 \(g\in G\)，使得 \(hBh^{-1}=gBg^{-1}\) 且 \(hNh^{-1}=gNg^{-1}\)。令 \(\hat{B}\)（分别为 \(\hat{N}\)）为 \(\tilde{G}\) 中 B（分别为 N）的正规化子；置 \(\Gamma=\hat{B}\cap\hat{N}\)、\(\hat{N}=\Gamma.N\) 和 \(\tilde{T}=\tilde{N}\cap B\)。
\end{enumerate}

\begin{enumerate}
\def\labelenumi{\alph{enumi})}
\item
  证明 \(\tilde{\mathbf{G}} = \Gamma .\mathbf{G}\)、\(\hat{\mathbf{B}} = \Gamma .\mathbf{B}\)、\(\Gamma \cap \mathrm{B} = \Gamma \cap \mathrm{G}\) 以及 \(\tilde{\mathbf{T}} = (\Gamma \cap \mathbf{B}).\mathbf{T}\)。因此，群 \(\Omega = \Gamma /(\Gamma \cap \mathrm{B})\)、\(\tilde{\mathbf{G}} /\mathbf{G}\) 和 \(\hat{\mathbf{B}} /\mathbf{B}\) 典范同构。若 \(\varPhi \subset \varOmega\)，且 \(\mathbf{H}\) 是 \(\mathbf{G}\) 的包含 \(\Gamma \cap \mathrm{B}\) 的子群，则以 \(\varPhi \mathrm{H}\) 表示子集 \(\varphi \mathrm{H}\)（其中 \(\varphi \in \varPhi\)）的并集。
\item
  证明 \(\tilde{T}\) 在 \(\tilde{N}\) 中正规（为证明 \(n\gamma n^{-1} \in \tilde{N}\) 对 \(n \in \tilde{N}\) 和 \(\gamma \in \Gamma \cap B\) 成立，使用习题 5 a)），并且 \(N \cap \tilde{T} = T\)、\(\Gamma \cap \tilde{T} = \Gamma \cap B\)。因此，将 N（分别为 \(\Gamma\)）嵌入 \(\tilde{N}\)，可将 W（分别为 \(\Omega\)）认作 \(\tilde{W} = \tilde{N}/\tilde{T}\) 的一个子群。证明 \(\Omega\) 正规化 S，且 \(\tilde{W}\) 是 \(\Omega\) 与 W 的半直积。
\item
  证明对所有 \(s \in S\) 和所有 \(u \in \tilde{W}\)，
\end{enumerate}

\[
\mathrm{BsBuB} \subset (\mathrm{BuB}) \cup (\mathrm{BsuB}).
\]

\begin{enumerate}
\def\labelenumi{\alph{enumi})}
\setcounter{enumi}{3}
\item
  证明映射 \(u \mapsto BuB\) 是从 \(\tilde{W}\) 到 \(B\backslash\tilde{G}/B\) 的双射（使用定理 1 以及 \(\Gamma\) 正规化 B 这一事实）。
\item
  令 \(\mathfrak{B}\) 为二元组 \((\varPhi, X)\) 的集合，其中 \(\varPhi\) 是 \(\Omega\) 的子群，且 \(X\) 是 \(X\) 的子集并被 \(\varPhi\) 正规化。置 \(G_{(\varPhi, X)} = \varPhi G_X = B \varPhi W_X B\)（沿用第 5 项的记号）。证明映射 \((\varPhi, X) \mapsto G_{(\varPhi, X)}\) 是从 \(\mathfrak{B}\) 到 \(\tilde{G}\) 中包含 \(B\) 的子群集合的双射。推广第 5 项定理 3 的断言 \(b)\) 和 \(c)\)，以及命题 2。
\end{enumerate}

证明 \(\mathrm{G}_{(\varPhi,\mathrm{X})}\) 在 \(\tilde{G}\) 中的正规化子是子群 \(\mathrm{G}_{(\varPhi',\mathrm{X})}\)，其中 \(\varPhi'\) 是 \(\Omega\) 中同时正规化 \(\varPhi\) 和 X 的元素所成的集合。

\begin{enumerate}
\def\labelenumi{\alph{enumi})}
\setcounter{enumi}{5}
\tightlist
\item
  证明 \(G_{(\Phi, \mathsf{X})}\) 是 \(\tilde{\mathbf{G}}\) 的极大子群，当且仅当以下两个条件之一得到满足：
\end{enumerate}

\begin{enumerate}
\def\labelenumi{(\roman{enumi})}
\item
  \(X = S\) 且 \(\varPhi\) 是 \(\varOmega\) 的极大子群；
\item
  \(\Phi = \Omega\) 且 \(\Phi\) 在 \(S - X\) 上传递地作用（该集合非空）。
\end{enumerate}

证明 \(G_{(\Phi,X)}\) 在不包含 G 的 \(\tilde{G}\) 的子群集合中是极大子群，当且仅当

\begin{enumerate}
\def\labelenumi{(\roman{enumi})}
\setcounter{enumi}{2}
\tightlist
\item
  \(X \neq S\)，\(\Phi\) 是 \(X\) 在 \(\Omega\) 中的正规化子，且在 \(S - X\) 上传递地作用。
\end{enumerate}

\(g)\) 设 \(\Phi\) 为 \(\Omega\) 的正规子群，并置 \(G' = \Phi G\)、\(B' = \Phi B\)、\(N' = \Phi N\) 和 \(T' = B' \cap N'\)。证明 \(T' = \Phi \tilde{T}\)，且 \(T'\) 在 \(N'\) 中正规，当且仅当 \(\Phi\) 的每个元素都与 \(W\) 的每个元素交换。证明此时 \(G'\) 在 \(\tilde{G}\) 中正规，且 \(N\) 到 \(N'\) 的嵌入定义了同构 \(j\)，它从 \(W\) 到 \(W' = N'/T'\)；并且 \((G', B', N', S')\)（其中 \(S' = j(S)\)）是 Tits 系统。

\begin{enumerate}
\def\labelenumi{\arabic{enumi})}
\setcounter{enumi}{8}
\tightlist
\item
  令 \((G, B, N, S)\) 为 Tits 系统，X、Y、Z 为 S 的三个子集。证明
\end{enumerate}

\[
\mathrm{G} _ {\mathrm{X}} \cap (\mathrm{G} _ {\mathrm{Y}}. \mathrm{G} _ {\mathrm{Z}}) = (\mathrm{G} _ {\mathrm{X}} \cap \mathrm{G} _ {\mathrm{Y}}). (\mathrm{G} _ {\mathrm{X}} \cap \mathrm{G} _ {\mathrm{Z}})
\]

（使用 §1 的习题 1 以及第 3 项的命题 2）。

\begin{enumerate}
\def\labelenumi{\arabic{enumi})}
\setcounter{enumi}{9}
\tightlist
\item
  设 G 是一个群，B 是 G 的 Tits 子群。沿用习题 3 的记号。对 \(s \in S\)，以 \(G^{(s)}\) 表示子群 \(G_{S - \{s\}}\)（习题 3 g)）。令 I 为形如 \(gG^{(s)}\)（其中 \(g \in G\)、\(s \in S\)）的 G 的子集的集合，令 C 为形如 \(C_g = \{gG^{(s)} \mid s \in S\}\)（其中 \(g \in G\)）的 I 的子集的集合。称 \(C_g\) 为 I 的室（§ 1，习题 15）。令 G 通过左平移作用于 I。
\end{enumerate}

\begin{enumerate}
\def\labelenumi{\alph{enumi})}
\item
  令 \(\mathfrak{F}\) 为 I 的面的集合（即室的子集，参见 §1，习题 15）。令 \(F \in \mathfrak{F}\)。证明存在 S 的唯一子集 X 和 \(g \in G\)，使得 \(gG_{X} = \bigcap_{a \in F} a\)；称 F 为 X 型。证明此时 F 正是 \(gG^{(s)}\)（其中 \(s \in S - X\)）的集合，并且它在包含它的任一室中的余维为 \(\operatorname{Card}(X)\)。证明映射 \(j: F \mapsto \bigcap_{a \in F} a\) 是从 \(\mathfrak{F}\) 到形如 \(gG_{X}\)（其中 \(g \in G\)、\(X \subset S\)）的 G 的子集的严格递减双射，并与左平移相容。证明：若 \(X \subset Y \subset S\)，则每个 X 型面都包含唯一的 Y 型面。
\item
  证明 G 在室集 C 上传递地作用，并且室 \(C_{g}\)（\(g \in G\)）的稳定子群等于 \(gBg^{-1}\)；因此映射 \(g \mapsto C_{g}\) 定义了从 G/B 到 C 的双射。
\item
  证明两个室 \(C_{g}\) 和 \(C_{g'}\)（\(g, g' \in G\)）相邻（§1，习题 15），当且仅当存在 \(s \in S\) 使得 \(g' \in g(B \cup BsB)\)。
\end{enumerate}

\(d)\) 令 \(C_1, \ldots, C_n\) 为 I 的室，并置 \(C_0 = C_e\)。证明以下条件等价：

\begin{enumerate}
\def\labelenumi{(\roman{enumi})}
\item
  序列 \(\Gamma = (\mathrm{C}_0, \mathrm{C}_1, \ldots, \mathrm{C}_n)\) 是一个单射廊；
\item
  存在由 S 的元素组成的序列 \(\mathbf{s} = (s_1, \ldots, s_n)\) 和由 B 的元素组成的序列 \((b_1, \ldots, b_n)\)，使得 \(C_j = b_1 \bar{s}_1 b_2 \bar{s}_2 \ldots b_j \bar{s}_j (C_0)\)（其中 \(\bar{s}_j\) 表示双陪集 \(Bs_jB\) 的一个给定元素）对 \(1 \leqslant j \leqslant n\) 成立。
\end{enumerate}

证明：若这些条件满足，则序列 s 是唯一的：称其为 \(\Gamma\) 的类型，记为 \(\mathbf{s}(\Gamma)\)。证明单射廊是最短的，当且仅当其类型是一个约化分解。证明以下条件得到满足：

(WI 1) 对 I 的任意两个室 C 和 \(\mathbf{C}'\)，存在唯一的 W 的元素 \(t(\mathbf{C},\mathbf{C}')\)，使得端点为 C 和 \(\mathbf{C}'\) 的最短廊的类型集正是 \(t(\mathbf{C},\mathbf{C}')\) 的约化分解集。

(WI 2) 对任意室 C，映射 \(\mathrm{C}' \mapsto t(\mathrm{C}, \mathrm{C}')\) 从 I 的室集到 W 是满射。

\begin{enumerate}
\def\labelenumi{\alph{enumi})}
\setcounter{enumi}{4}
\item
  证明两个类型相同且端点相同的最短廊相同。（归结为证明：若 \((s_{1},\ldots,s_{n})\) 是一个约化分解，且 \(b_{1},\ldots,b_{n},b_{1}^{\prime},\ldots,b_{n}^{\prime}\) 是 B 的元素，满足 \(b_{1}\bar{s}_{1}\ldots b_{n}\bar{s}_{n}\in b_{1}^{\prime}\bar{s}_{1}\ldots b_{n}^{\prime}\bar{s}_{n}B\)，则 \(b_{1}\bar{s}_{1}\in b_{1}^{\prime}\bar{s}_{1}B\)。为此，注意若 \(\bar{s}_{1}^{-1}b_{1}^{-1}b_{1}^{\prime}\bar{s}_{1}\notin B\)，则此元素将属于 \(Bs_{1}B\)，于是有 \(b_{2}\bar{s}_{2}\ldots b_{n}\bar{s}_{n}\in b\bar{s}_{1}b^{\prime}b_{2}^{\prime}\bar{s}_{2}\ldots b_{n}^{\prime}\bar{s}_{n}B\)，其中 \(b,b^{\prime}\in B\)，这与第 4 项定理 2 的推论 1 矛盾。）
\item
  证明 I 配备 C 后是一个建筑，称为与 (G, B) 相关的建筑。证明 I 存在唯一的编号（§ 1，习题 20），使得一个面的类型就是 a) 中定义的类型。
\item
  证明 I 是宽敞的，即 I 的每个面板至少包含于三个室中（参见 § 1，习题 24）。证明以下条件得到满足：
\end{enumerate}

\begin{enumerate}
\def\labelenumi{(\Alph{enumi})}
\setcounter{enumi}{6}
\tightlist
\item
  给定一个面板 F、一个室 C、一个廊 \(\Gamma = (\mathrm{C}_0, \ldots, \mathrm{C}_n)\)，满足 \(\mathrm{F} \subset \mathrm{C}_n\) 且长度最短，以及包含 F 的室 \(\mathrm{C}'\) 和 \(\mathrm{C}''\)，它们都异于 \(\mathrm{C}_n\)，存在元素 \(g \in \mathrm{G}\)，使得 \(g(\mathrm{C}_i) = \mathrm{C}_i\) 对 \(0 \leqslant i \leqslant n\) 成立，且 \(g(\mathrm{C}') = \mathrm{C}''\)。
\end{enumerate}

(归结为 \(C_{0} = C\) 的情形；令 \(u \in S\) 使 F 为 \(S - \{u\}\) 型，并令 s 为 \(\Gamma\) 的类型。证明 \((\mathbf{s}, u)\) 是一个约化分解。取 \(h \in G\) 使得 \(C_{n} = h(\mathrm{C})\)；存在 \(b'\) 和 \(b'' \in B\)，使得 \(C' = hb'u(\mathrm{C})\) 且 \(C'' = hb''u(\mathrm{C})\)。然后使用推广到 Tits 子群情形的第 4 项定理 2 的推论 1（参见习题 3 f)），证明存在 \(b \in B\) 使得 \(bhb'uB = hb''uB\)；于是 \(b \in hBuBu^{-1}Bh^{-1} \subset (hBh^{-1}) \cup (hBuBh^{-1})\)。若 \(b \in hBuBh^{-1}\)，则有 BhB \(\subset BhBuB = BhuB\)（同上引），这是不可能的。因此 \(b(C_{n}) = C_{n}\)，由 e) 得对所有 i 都有 \(b(C_{i}) = C_{i}\)。）

\begin{enumerate}
\def\labelenumi{\arabic{enumi})}
\setcounter{enumi}{10}
\tightlist
\item
  设 \((\mathbf{W},\mathbf{S})\) 为 Coxeter 系统，I 为由 S 编号的建筑（§ 1，习题 20）。若 \(\Gamma = (\mathbf{C}_0,\dots ,\mathbf{C}_n)\) 是一个单射廊，则 \(\Gamma\) 的类型（记为 \(s(\Gamma)\)）是 S 的元素组成的序列 \((s_1,\ldots ,s_n)\)，满足面板 \(\mathbf{C}_{i - 1}\cap \mathbf{C}_i\) 的类型为 \(\mathbf{S} - \{s_i\}\)（当 \(1\leqslant i\leqslant n\) 时）。若习题 10 的条件 (WI 1) 和 (WI 2) 得到满足，则称 I 为一个 \((\mathbf{W},\mathbf{S})\)-建筑。
\end{enumerate}

设 I 是一个宽敞的 (W, S)-建筑（参见习题 10 g)），G 是 I 的满足以下条件的允许自同构群：

(G0) 给定包含同一面板的三个不同室 C、C' 和 C''，存在元素 \(g \in G\)，使得 \(g(\mathrm{C}) = \mathrm{C}\) 且 \(g(\mathrm{C}') = \mathrm{C}'\)。

取 I 的一个室 C，并以 B 表示 C 在 G 中的稳定子群。

\begin{enumerate}
\def\labelenumi{\alph{enumi})}
\item
  证明 G 在 I 的室集上传递地作用。
\item
  设 \(C'\) 和 \(C''\) 是两个室。证明 \(t(\mathrm{C},\mathrm{C}') = t(\mathrm{C},\mathrm{C}'')\)，当且仅当存在 \(b \in B\) 使得 \(\mathrm{C}'' = b(\mathrm{C}')\)（若 \(t(\mathrm{C},\mathrm{C}') = t(\mathrm{C},\mathrm{C}'') = w\)，取 w 的一个约化分解 s，以及类型为 s 的最短廊 \(\Gamma'\) 和 \(\Gamma''\)，端点分别为 C 与 \(C'\)、C 与 \(C''\)。利用 \(l(w)\) 作归纳，使用 \((\mathrm{G}_{0})\) 和 a)。推出存在双射 \(w \mapsto \mathrm{B}(w)\)，从 W 到 B\textbackslash G/B。
\item
  令 F 为 C 的 S- \(\{s\}\) 型面。证明 F 在 G 中的稳定子群等于 B \(\cup\) B(s)（若 g(F) = F，则 t(C, g(C)) = 1 或 s）。证明 B \(\subset\) B(s)B(s)（使用 (G0)）。
\item
  证明 B 是 G 的 Tits 子群，且 (G,B) 的 Coxeter 系统典范同构于 (W,S)。（令 \(w \in W\) 和 \(s \in S\) 满足 \(l_{\mathbb{S}}(sw) = l_{\mathbb{S}}(w) + 1\)。令 \(g \in \mathrm{B}(w)\)、\(u \in \mathrm{B}(s)\)；令 \(\mathbf{s} = (s_1, \ldots, s_n)\) 为 \(w\) 的一个约化分解，\(C = C_0, C_1, \ldots, C_n = g(C)\) 为类型为 s 的最短廊，并取 \(\bar{s}_i \in \mathrm{B}(s_i)\)；证明利用 \((G_0)\)，存在 \(b_i \in B\)，使得 \(C_i = b_1 \bar{s}_1 \ldots b_i \bar{s}_i(C)\)。置 \(C_i' = u(C_{i-1})\)，证明廊 \((C, u(C), u(C_1), \ldots, u(C_n))\) 的类型为 \((s, s)\)，因而是最短的。推出 \(ug \in \mathrm{B}(sw)\) 且 \(B(s) B(w) = B(sw)\)。若现在 \(l(sw) = l(w) - 1\)，置 \(w' = sw\)：于是 \(B(s) B(w') = B(w)\)，从而
\end{enumerate}

\[
\mathrm{B} (s) \mathrm{B} (w) = \mathrm{B} (s) \mathrm{B} (s) \mathrm{B} \left(w ^ {\prime}\right) \subset \mathrm{B} \left(w ^ {\prime}\right) \cup \left(\mathrm{B} (s) \mathrm{B} \left(w ^ {\prime}\right)\right) = \mathrm{B} (s w) \cup \mathrm{B} (w);
\]

最后，由于 \(\mathrm{B} \subset \mathrm{B}(s)\mathrm{B}(s)\)，还有 \(\mathrm{B}(w') = \mathrm{B}(sw) \subset \mathrm{B}(s)\mathrm{B}(w)\)。

\begin{enumerate}
\def\labelenumi{\alph{enumi})}
\setcounter{enumi}{4}
\tightlist
\item
  证明存在从与 (G, B) 相关的建筑（习题 10）到 I 的唯一同构，它与 G 的作用相容，并将典范室 \(C_{e}\) 映到 C。
\end{enumerate}

\begin{enumerate}
\def\labelenumi{\arabic{enumi})}
\setcounter{enumi}{11}
\tightlist
\item
  设 \((\mathbf{G},\mathbf{B},\mathbf{N},\mathbf{S})\) 为一个 Tits 系统，\(\mathrm{W} = \mathrm{N} / (\mathrm{B}\cap \mathrm{N})\) 为其 Weyl 群，I 为与 \((\mathrm{W},\mathrm{S})\) 相关的 \((\mathbf{G},\mathbf{B})\)-建筑（习题 10），\(\mathbf{C} = \mathbf{C}_e\) 为 I 的典范室。令 \(\mathbf{A}_0\) 为与 Coxeter 系统 \((\mathrm{W},\mathrm{S})\) 相关的公寓（§ 1，习题 16）。对所有 \(g\in \mathbf{G}\)，令 \(\varphi_g\) 为从 \(\mathbf{A}_0\) 到 I 的映射；它将点 \(w\mathrm{W}^{(s)}\)（属于 \(\mathbf{A}_0\)，其中 \(w\in \mathrm{W}, s\in \mathrm{S}\)）映到 I 的点 \(gw\mathrm{G}^{(s)}\)。
\end{enumerate}

证明：对所有 \(g \in G\)，映射 \(\varphi_{g}\) 是从 \(A_{0}\) 到 I 的一个有编号建筑的同构，其像是由 \(gn(C_{e})\)（其中 \(n \in N\)）组成的室的并集。证明 I 配备由 \(\varphi_{g}(A_{0})\)（其中 \(g \in G\)）组成的集合 U 后，是一个结构建筑（§1，习题 24）（为证明 (SB 2)，注意若 \(g', g'' \in G\)，则存在 \(b', b'' \in B\) 和 \(n \in N\)，使得 \(g'^{-1}g'' = b'nb''\)；于是置 \(g = g'b'n\)，并证明 \(g'(C)\) 和 \(g''(C)\) 包含于 \(\varphi_{g}(A_{0})\) 中。为证明 (SB 3)，归结到两个公寓 \(A' = \varphi_{e}(A_{0})\) 和

\(A'' = \varphi_b(A_0)\)，其中 \(b \in B\)，并使用第 5 项命题 2 证明映射 \(a \mapsto b(a)\) 固定 \(A' \cap A''\) 的各点）。证明 Coxeter 系统 (W, S) 和 I 的编号适配于 (I, U)（§1，习题 24 e)），并且集合 \(\mathfrak{I}\)，即从 \(A_0\) 到 \(\mathfrak{U}\) 的各个元素的允许同构，正是 \(\varphi_g\)（其中 \(g \in G\)）的集合。

\begin{enumerate}
\def\labelenumi{\arabic{enumi})}
\setcounter{enumi}{12}
\tightlist
\item
  沿用 §1 习题 24 的记号：(I, U) 是一个配备 Coxeter 系统 (W, S) 和适配编号的宽敞结构建筑，而 I 是从与 (W, S) 相关的公寓 \(A_{0}\) 到 (I, U) 的各个公寓的允许同构的集合。此外，令 G 为 I 的保持 E 的允许自同构群：于是 G 作用于 I，并假设 G 在 I 上传递地作用。
\end{enumerate}

以 C 表示 I 的一个室，A 表示 U 中包含 C 的一个公寓，以 \(\varphi\) 表示将 \(A_{0}\) 的典范室 \(C_{e}\) 映到 C 的允许同构（从 \(A_{0}\) 到 A）；以 B 表示 C 在 G 中的稳定子群，以 N 表示 A 在 G 中的稳定子群。

\begin{enumerate}
\def\labelenumi{\alph{enumi})}
\item
  证明：若 \(A'\) 和 \(A''\) 是 U 中包含同一室的两个公寓，则存在 \(g \in G\)，使得 \(g(A') = A''\)，且 \(g(a) = a\) 对所有 \(a \in A' \cap A''\) 成立。证明 G 在二元组 (A, C) 的集合上传递地作用，其中 \(A \in U\)，C 是 A 的一个室。
\item
  证明映射 \(n \mapsto \varphi^{-1} \circ n \circ \varphi\) 是从 N 到 W 的满同态，其核为 \(B \cap N\)。因此可以将 \(N/(B \cap N)\) 认作 W。
\item
  证明习题 10 d) 的条件 (WI 1) 和 (WI 2) 得到满足（使用以下事实：若 I 的一个公寓包含两个室 C' 和 C''，则它包含端点为 C' 和 C'' 的每个最短廊（§ 1，习题 24 b)）。
\item
  证明习题 10 g) 的条件 (G)（从而条件 (G₀) 也得到满足）得到满足（沿用 (G) 的记号，考虑 I 中包含 C₀ 和 C'（分别为 C''）的公寓 A'（分别为 A''）；使用 §1 习题 24 b) 证明 Γ ⊂ A'' ∩ A''，并使用 a)）。
\item
  证明 \((G, B, N, S)\) 是一个 Tits 系统，且 \((I, \mathfrak{U})\) 典范同构于相关的有编号结构建筑（习题 12）。
\end{enumerate}

\begin{enumerate}
\def\labelenumi{\arabic{enumi})}
\setcounter{enumi}{13}
\item
  设 \((\mathbf{G},\mathbf{B},\mathbf{N},\mathbf{S})\) 为一个 Tits 系统，\((\mathbf{I},\mathfrak{U})\) 为相关的有编号结构建筑（习题 12）。证明 G 视为 I 的自同构群时满足习题 13 的条件。沿用习题 12 的记号，置 \(\mathrm{A}=\varphi_{e}(\mathrm{A}_{0})\) 和 \(\mathrm{C}=\varphi_{e}(\mathrm{C}_{e})\)；证明 B 是 C 在 G 中的稳定子群。令 \(\tilde{N}\) 为 A 在 G 中的稳定子群；证明 \(\tilde{\mathrm{N}}/(\mathrm{B}\cap\tilde{\mathrm{N}})\) 可以认作 W，且 \((\mathrm{G},\mathrm{B},\tilde{\mathrm{N}},\mathrm{S})\) 是与 \((\mathrm{G},\mathrm{B},\mathrm{N},\mathrm{S})\) 相关的饱和 Tits 系统（习题 5）。
\item
  沿用习题 10 的假设和记号。此外，假设 \(\mathbf{W}\)（\((\mathbf{G},\mathbf{B})\) 的 Weyl 群）有限，并以 \(w_0\) 表示 \(\mathbf{W}\) 的最长元素（§1，习题 22）。若两个室 \(\mathbf{C}\) 和 \(\mathbf{C}'\) 满足 \(t(\mathbf{C},\mathbf{C}') = w_0\)，则称其对置。
\end{enumerate}

\begin{enumerate}
\def\labelenumi{\alph{enumi})}
\item
  证明若 C 与 C' 对置，则 C' 与 C 也对置。证明每个给定的室 C 都存在一个与之对置的室。证明室 C 在 G 中的稳定子群在与 C 对置的室集上是传递的。
\item
  设 C 与 \(C'\) 是两个对置的室。证明对所有 \(w \in W\)，存在唯一的室 \(C_w\) 具有如下性质：若 \((s_1, \ldots, s_k)\)（分别为 \((s'_1, \ldots, s'_h)\)）是 w（分别是 \(w' = w_0 w^{-1}\)）的一个约化分解，则存在最短廊 \((C_0 = C, C_1, \ldots, C_n = C')\)，其类型为 \((s_1, \ldots, s_k, s'_1, \ldots, s'_h)\)（其中 \(n = h + k\)），使得 \(C_w = C_k\)（使用 §1 习题 22 以及习题 10 d) 和 e)）。证明 \(C_w\) 与 \(C_{ww_0}\) 对置。
\item
  令 \(\mathfrak{M}\) 为对置室的二元组集合。对 \(m = (C, C') \in \mathfrak{M}\)，令 \(A_m\) 为上面构造的室 \(C_w\) 的并集。证明 I 配备由 \(\mathfrak{U}\) 中的 \(A_m\)（其中 \(m \in \mathfrak{M}\)）组成的集合后是一个结构建筑（\(\S 1\)，习题 24），且 Coxeter 系统 (W, S) 和 I 的编号适配于 (I, \(\mathfrak{U}\) )；定义从 \(\mathfrak{M}\) 到习题 24 中记为 \(\mathfrak{I}\) 的集合（该习题见 \(\S 1\)）的典范双射。我们认同这两个集合。
\item
  令 \(m = (C, C') \in \mathfrak{M}\)，其中 \(C = C_e\)，并令 N 为 \(A_m\) 在 G 中的稳定子群。证明 \(N/(B \cap N)\) 可以认作 W，且 \((G, B, N, S)\) 是饱和的 Tits 系统（使用习题 13）。
\end{enumerate}

\begin{enumerate}
\def\labelenumi{\arabic{enumi})}
\setcounter{enumi}{15}
\tightlist
\item
  沿用习题 15 的假设和记号。
\end{enumerate}

\begin{enumerate}
\def\labelenumi{\alph{enumi})}
\item
  设 C 和 \(C'\) 是两个室。证明存在一个室 \(C''\)，同时与 C 和 \(C'\) 对置（取与 C 对置的 \(C''\)，使 \(t(\mathrm{C},\mathrm{C}')\) 具有可能的最大长度；若 \(t(\mathrm{C}',\mathrm{C}'')\neq w_{0}\)，则存在一个室 \(C_{1}\)，它与 \(C''\) 相邻，使得 \(l(t(\mathrm{C}',\mathrm{C}_{1}))>l(t(\mathrm{C}',\mathrm{C}''))\)；使用习题 10 的条件 (G) 证明可以假设 \(C_{1}\not\subset A_{(C,C'')}\)，并且此时 \(C_{1}\) 与 C 对置）。
\item
  令 \(a \in U\)；证明存在唯一的对合自同构 \(j_{A}\)（不一定允许），将 A 的每个室映到其对置室（使用 §1 习题 22 c)）。令 F 和 \(F'\) 为 I 的两个面；证明若 \(F' = j_{A}(F)\) 对某个 \(A \in U\) 成立且该 A 包含 F 和 \(F'\)，则对所有 \(A \in U\)（包含 F 和 \(F'\)）都成立同样结论（若 \(F, F' \in A \cap A'\)，其中 \(A, A' \in U\)，则考虑室 C（分别为 \(C'\)）（分别属于 A 和 \(A'\)，并分别包含 F 和 \(F'\)），再取 \(A'' \in U\) 包含 C 和 \(C'\)，并使用 (SB 3)）。于是称 F 与 \(F'\) 对置。证明两个面有公共的对置面，当且仅当它们具有相同的类型 T；并且与 T 型面相对的面是 \(w_{0}Tw_{0}^{-1}\) 型。
\item
  令 \(A_{0}\) 为与 Coxeter 系统 (W, S) 相关的公寓（§1，习题 16），令 J 为从 \(A_{0}\) 到 U 的各个元素的允许同构的集合。若 \(\alpha\) 是 \(A_{0}\) 的一个半部，其墙为 L（§1，习题 16），且 \(\varphi \in J\)，则称 \(\varphi(\alpha)\) 为 I 的一个半公寓，其墙为 \(\varphi(L)\)。证明 \(\varphi(L)\) 只依赖于 \(\varphi(\alpha)\)，而不依赖于二元组 \((\varphi, \alpha)\)。令 \(D_{1}\) 和 \(D_{2}\) 为具有同一墙 L 的两个不同半公寓：证明存在 \(\varphi \in J\) 和 \(L_{0}\)（属于 \(A_{0}\) 的一面墙），使得 \(L = \varphi(L_{0})\) 并且 \(D_{i}\) 是 \(\varphi\) 将 \(A_{0}\) 中由 \(L_{0}\) 确定的两个半部映出的像（取包含于 L 中的一个面板，以及室 \(C_{1}\) 和 \(C_{2}\)，其中一个包含于 \(D_{1}\)，另一个包含于 \(D_{2}\)，并且 \(C_{1}\) 包含 F，\(C_{2}\) 包含 \(D_{1}\) 中与 F 对置的面板（该面板在 \(D_{2}\) 中也与 F 对置），再考虑包含 \(C_{1}\) 和 \(C_{2}\) 的 I 的公寓）。
\end{enumerate}

\begin{enumerate}
\def\labelenumi{\arabic{enumi})}
\setcounter{enumi}{16}
\tightlist
\item
  沿用习题 15 和 16 的假设及记号。取 \(\varphi \in \mathfrak{I}\)，将典范室 \(C\)（属于 \(A_0\)，见 \(\S 1\)，习题 16）映到典范室 \(C_e\)（属于 \(I\)），并如习题 15 \(d\)）中那样，以 \(N\) 表示公寓 \(\varphi(A_0) \in \mathfrak{U}\) 在群 \(G\) 中的稳定子群。对任意子集 \(D\)（属于 \(A_0\)），以 \(B_D\) 表示 \(G\) 中固定 \(\varphi(D)\) 的所有点的子群：有 \(B = B_C\) 且 \(B \cap N = B_{A_0}\)。
\end{enumerate}

\begin{enumerate}
\def\labelenumi{\alph{enumi})}
\item
  设 \(\alpha\) 是包含 C 的 \(A_0\) 的一个半部。证明 \(B_{\alpha}\) 在 I 中与 \(\varphi(\alpha)\) 不同且其墙为 \(\varphi(\alpha)\) 的墙 L 的半公寓集合上传递地作用（令 X 为这样的半公寓，F 为包含于 L 中的面板，\(C'\) 为包含 F 的 \(\varphi(\alpha)\) 的一个室；使用习题 16 c) 和习题 13 a) 证明存在 \(g \in G\)，使得 \(g(X \cup \varphi(\alpha)) = \varphi(A_0)\) 且 \(g(C') = C'\)；证明 \(g(F) = F\)，并由 §1 习题 22 c) 推出 \(g \in B_{\alpha}\)。推出 \(B_L\) 在墙为 L 的半公寓上二重传递地作用，且 \((B_L, B_{\alpha}, B_L \cap N)\) 是 Weyl 群阶为 2 的 Tits 系统（参见习题 7）。
\item
  设 \(D_{1}\) 和 \(D_{2}\) 是 \(A_{0}\) 的两个凸子集，满足 \(C \subset D_{1} \subset D_{2}\)，并且存在唯一的半部 \(\alpha\)（属于 \(A_{0}\)），使得 \(D_{1} \subset \alpha\) 而 \(D_{2} \not\subset \alpha\)：于是 \(D_{1} = D_{2} \cap \alpha\)（§1，习题 21 c)）。证明 \(B_{D_{1}} = B_{\alpha}B_{D_{2}}\)，且 \(B_{\alpha} \cap B_{D_{2}} = B \cap N\)（考虑一个最短廊，其一个端点等于 C，另一个端点包含点 \(a \in D_{2} - D_{1}\)，证明存在两个相邻室 \(C_{1}'\) 和 \(C_{2}'\)（属于 \(A_{0}\)），满足 \(C_{1}' \subset D_{1}\)、\(C_{2}' \subset D_{2}\)、\(C_{2}' \not\subset D_{1}\)，且 \(C_{1}' \cap C_{2}'\) 包含于墙 \(L'\)（属于 \(\alpha\)）中。置 \(C_{i} = \varphi(C_{i}')\)、\(F = \varphi(C_{1}' \cap C_{2}')\) 和 \(L = \varphi(L')\)。证明 \(D_{1} \cup C_{2}'\) 的凸包等于 \(D_{2}\)。令 \(b \in B_{D_{1}}\)：则 \(b(C_{1}) = C_{1}\)，因而 \(b(F) = F\) 且 \(b(C_{2}) \neq C_{1}\)。推出 \(b(C_{2}) \cup L\) 的凸包是异于 \(\varphi(\alpha)\) 且墙为 L 的半公寓 X，并且存在 \(b' \in B_{\alpha}\)，使得
\end{enumerate}

\[
b ^ {\prime} (\varphi (A _ {0})) = X \cup \varphi (\alpha).
\]

证明 \(b'b(a) = a\) 对所有 \(a \in \varphi(\mathrm{D}_1 \cup \mathrm{C}_2')\) 成立，且 \(b'b \in \mathrm{B}_{\mathrm{D}_2}\) 。

\begin{enumerate}
\def\labelenumi{\alph{enumi})}
\setcounter{enumi}{2}
\tightlist
\item
  令 \((\alpha_{i})_{1\leqslant i\leqslant q}\) 为包含 C 的 \(A_0\) 的半部，并对其编号，使得映射
\end{enumerate}

\[
j \mapsto \bigcap_ {1 \leqslant i \leqslant j} \alpha_ {i}
\]

  \(\mathbf{B} = \mathbf{B}_{\alpha_1} \ldots \mathbf{B}_{\alpha_q}\)。证明若 \(b_i, b_i' \in \mathbf{B}_{\alpha_i}\) 且 \(b_1 \ldots b_q = b_1' \ldots b_q'\)，则 \(b_i' \in b_i(\mathbf{B} \cap \mathbf{N})\) 对 \(1 \leqslant i \leqslant q\) 成立。

\begin{enumerate}
\def\labelenumi{(\arabic{enumi})}
\setcounter{enumi}{17}
\tightlist
\item
  令 \(V_i\) 为由 \(k^n\) 的基 \(e_1, \ldots, e_i\) 张成的子空间（\(1 \leqslant i \leqslant n - 1\)）。
\end{enumerate}

\begin{enumerate}
\def\labelenumi{\alph{enumi})}
\item
  证明 \(G_{X}\) 是满足 \(g \in G\) 且 \(g(V_{i}) = V_{i}\) 时有 \(s_{i} \notin X\) 的元素所成的集合。
\item
  令 I 为与 Tits 系统 (G, B, N, S) 相关的建筑（习题 10 和 12）。证明将 I 的点 \(g\mathbf{G}^{(i)}\)（其中 \(\mathbf{G}^{(i)}\) 表示 G 的子群 \(G_{S-\{s_i\}}\)，换言之是 \(V_i\) 的稳定子群）关联到向量子空间 \(g(V_i)\) 的映射 j，是从 I 到 \(k^n\) 的满足 \(\neq \{0\}\) 且 \(\neq k^n\) 的向量子空间集合的双射，并与 G 的作用相容。
\item
  设 E 为向量空间。E 的一个旗是 E 的向量子空间的集合，且按包含关系全序。证明 I 的元素 \(a_{1}, \ldots, a_{k}\) 属于 I 的同一个面，当且仅当 \(\{j(a_{1}), \ldots, j(a_{k})\}\) 是 \(k^{n}\) 的一个旗。
\item
  证明 G 在 \(k^{n}\) 的一维向量子空间集上二重传递地作用。假设 \(n \geqslant 2\)，令 \(N_{1}\) 为由交换 \(e_{1}\) 和 \(e_{2}\) 且固定其余 \(e_{i}\) 的元素生成的 G 的子群。证明 \((\mathrm{G}, \mathrm{G}^{(1)}, \mathrm{N}_{1})\) 是 Weyl 群阶为 2 的 Tits 系统。
\item
  假设 \(k\) 可交换。置
\end{enumerate}

\[
\mathrm{G} ^ {\prime} = \mathbf {S L} (n, k), \quad \mathrm{B} ^ {\prime} = \mathrm{G} ^ {\prime} \cap \mathrm{B}, \quad \mathrm{N} ^ {\prime} = \mathrm{G} ^ {\prime} \cap \mathrm{N} \text {and} \mathrm{T} ^ {\prime} = \mathrm{N} ^ {\prime} \cap \mathrm{B} ^ {\prime} = \mathrm{T} \cap \mathrm{G} ^ {\prime}.
\]

证明 \(N^{\prime}/T^{\prime}\) 可以认作 \(S_{n}\)，且 \((G^{\prime}, B^{\prime}, N^{\prime}, S)\) 是一个 Tits 系统（按第 2 项的方法论证）。

\begin{enumerate}
\def\labelenumi{\arabic{enumi})}
\setcounter{enumi}{18}
\item
  令 \(k\) 为特征 \(\neq 2\) 的可交换域，令 \(\mathbf{Q}\) 为二次型 \(x_{1}x_{3} + x_{2}^{2}\)，定义在 \(k^{3}\) 上。证明群 \(\mathbf{SO}(\mathbf{Q})\) 在 \(k^{3}\) 的各向同性直线集上二重传递地作用。将相应的 Tits 系统（习题 7）与第 2 项中 \(n = 2\) 时得到的 Tits 系统进行比较（参见《代数》第 IX 章，§9，习题 15）。
\item
  沿用第 2 项的记号，并假设 k 可交换。考虑以下情形：
\end{enumerate}

\((B_{l})\ n=2l+1\ (with\ l\geqslant1),\ k\ is\ of\ characteristic\neq2\ and\ k^{n}\ is\ equipped\) 配备二次型 \(Q=x_{1}x_{n}+\cdots+x_{l}x_{l+2}+x_{l+1}^{2};\)

\((C_{l}) \ n = 2l \ (\text{其中 } l \geqslant 1) \ \text{且} \ k^{n} \ \text{配备交错型 } \Phi \ \text{使得 } \Phi(e_{i}, e_{j}) = 0 \ \text{对} \ 1 \leqslant i \leqslant n, \ 1 \leqslant j \leqslant n \ \text{且} \ i \leqslant j, \ \text{除非} \ i \leqslant l, \ j = n + 1 - i, \ \text{此时} \Phi(e_{i}, e_{j}) = 1;\)

\((D_{l})\ n = 2l \text{ (其中 } l \geqslant 2), k \text{ 的特征 } \neq 2 \text{ 且 } k^{n} \text{ 配备二次型 } Q = x_{1}x_{n} + \cdots + x_{l}x_{l+1}.\)

以 \(G_{1}\) 表示特殊正交群 \(\mathbf{SO}(Q)\)（情形 \((B_{l})\) 和 \((D_{l})\)），并表示辛群 \(\mathbf{Sp}(\varPhi)\)（情形 \((C_{l})\)）。置

\[
\mathrm{B} _ {1} = \mathrm{G} _ {1} \cap \mathrm{B}, \quad \mathrm{N} _ {1} = \mathrm{G} _ {1} \cap \mathrm{N} \text {and} \mathrm{T} _ {1} = \mathrm{G} _ {1} \cap \mathrm{T} = \mathrm{B} _ {1} \cap \mathrm{N} _ {1}.
\]

\begin{enumerate}
\def\labelenumi{\alph{enumi})}
\tightlist
\item
  证明 \(N_{1}\) 在直线 \(ke_{i}\) 集合上的作用定义了一个同态：其定义域为 \(N_{1}\)，核为 \(T_{1}\)，且满射到 \(W_{1}\)（它是 \(S_{n}\) 的子群），从而可将 \(W_{1}\) 认作 \(N_{1}/T_{1}\)。证明 \(W_{1}\) 是由下列元素生成的 \(S_{n}\) 的子群：
\end{enumerate}

\(\sigma_{j} = s_{j}s_{n - j}\)（\(1\leqslant j <   l\)）以及 \(\sigma_l = s_{l - 1}s_ls_{l - 1}\)，情形 \((\mathrm{B}_l)\)；

\(\sigma_{j} = s_{j}s_{n - j}\)（\(1\leqslant j <   l\)）以及 \(\sigma_l = s_l\)，情形 \((C_l)\)；

\[
\text { the } \sigma_ {j} = s _ {j} s _ {n - j} \text { for } 1 \leqslant j <   l \text { and by } \sigma_ {l} = s _ {l - 1} s _ {l} s _ {l - 1} s _ {l} s _ {l + 1} s _ {l} \text { in case } (D _ {l}).
\]

\begin{enumerate}
\def\labelenumi{\alph{enumi})}
\setcounter{enumi}{1}
\tightlist
\item
  令 \(S_{1}\) 为由 \(\sigma_{j}\)（其中 \(1 \leqslant j \leqslant l\)）组成的集合。证明 \((\mathrm{G}_{1}, \mathrm{B}_{1}, \mathrm{N}_{1}, \mathrm{S}_{1})\) 是一个 Tits 系统（要证明 \(G_{1}\) 的由 \(B_{1}\) 和 \(N_{1}\) 生成的子群 H 等于整个 \(G_{1}\)，按《代数》第 II 章 §10 第 13 项论证），注意到 H 包含 \(G_{1}\) 的下三角子群，并且对于任意 \(\xi_{i} \in k\)（\(2 \leqslant i \leqslant n\)），存在矩阵 \(b = (b_{ij}) \in B_{1}\)，满足 \(b_{11} = 1\)、\(b_{1i} = \xi_{i}\)（\(2 \leqslant i \leqslant n - 1\)），且 \(b_{1n} = \xi_{n}\)（情形 \((\mathrm{C}_{l})\)），而 \(b_{1n} = 0\)（情形 \((\mathrm{B}_{l})\) 和 \((\mathrm{D}_{l})\)）。然后仿照第 2 项，引入子群 \(G_{1,j} = G_{1} \cap (G_{j}, G_{n-j})\)（\(1 \leqslant j < l\)）以及子群 \(G_{1,l}\)，其元素固定 \(G_{1}\)
\end{enumerate}

\(e_i\)（\(i \neq l, l + 1, l + 2\)）以及由 \(e_l, e_{l+1}\) 和 \(e_{l+2}\) 生成的子空间，情形 \((B_l)\)；

\(e_{i}\)（\(i \neq l, l + 1\)）以及由 \(e_{l}\) 和 \(e_{l+1}\) 生成的平面，情形 \((\mathrm{C}_{l})\)；\(e_{i}\)（对于 i \textless{} l - 1 或 \(i > l + 2\)）以及由 \(e_{l-1}\) 和 \(e_{l+1}\)、\(e_{l}\) 和 \(e_{l+2}\) 生成的两个平面，情形 \((\mathrm{D}_{l})\)。 

证明 \(G_{1,j}\) 分别可以认作 \(\mathbf{GL}(2,k)\)、\(\mathbf{SL}(2,k)\) 或习题 19 中的特殊正交群。

\(^{*}c)\) 证明群 \(W_{1}\) 的 Coxeter 图在上述三种情形中分别为 \((B_{l})\)、\((C_{l})\) 或 \((D_{l})\) 型（第 VI 章，§4，第 1 项）。

\(d\)）证明：对 \(\mathbf{S}_1\) 的任意子集 X，子群 \(\mathbf{G}_{1\mathbf{X}}\) 由满足 \(g\in \mathbf{G}_1\) 且 \(g(\mathrm{V}_i) = \mathrm{V}_i\) 对所有 \(i\)（其中 \(\sigma_{i}\notin \mathbf{X}\)）的元素组成；但在情形 \((\mathrm{D}_l)\) 中，若 \(\mathrm{V}_{r - 1}\) 表示由 \(e_1,\dots ,e_{r - 1}\) 和 \(e_{r + 1}\) 生成的子空间，则同一断言成立。仿照习题 18 b)，推出存在一个双射 \(j\)，从与 \((\mathrm{G}_1,B_1)\) 相关的建筑到满足 \(\neq 0\) 的全各向同性子空间集合（情形 \((\mathrm{B}_l)\) 和 \((\mathrm{C}_l)\)），以及维数满足 \(\neq 0\) 且 \(\neq r - 1\) 的全各向同性子空间集合（情形 \((\mathrm{D}_l)\)）。证明 I 的点 \(a_1,\ldots ,a_k\) 属于同一个面，当且仅当 \(\{j(a_1),\dots ,j(a_k)\}\) 是一个旗。

\begin{enumerate}
\def\labelenumi{\arabic{enumi})}
\setcounter{enumi}{20}
\tightlist
\item
  令 A 为离散赋值环（《交换代数》第 VI 章，§3，第 6 项），m 为其极大理想，γ 为 m 的一个生成元，K 为 A 的分式域。令 G 为群 SL(2,K)，B 为 G 的子群，由矩阵 \(\begin{pmatrix} a & b \\ c & d \end{pmatrix}\) 组成，其中 \(a, b, d \in A\) 且 \(c \in m\)（其中 ad - bc = 1）；令 N 为由属于 G 且每行每列恰有一个非零元素的矩阵组成的子群。
\end{enumerate}

\begin{enumerate}
\def\labelenumi{\alph{enumi})}
\item
  证明 \(\mathrm{T} = \mathrm{B} \cap \mathrm{N}\) 在 \(\mathbb{N}\) 中正规，且 \(\mathrm{W} = \mathrm{N} / \mathrm{T}\) 是由矩阵的类 \(s\) 和 \(s'\)（分别为 \(\left( \begin{array}{cc}0 & 1\\ -1 & 0 \end{array} \right)\) 和 \(\left( \begin{array}{cc}0 & \gamma \\ -\gamma^{-1} & 0 \end{array} \right)\)）生成的无限二面体群。
\item
  证明 \((G, B, N, S)\)（其中 \(S = \{s, s'\}\)）是一个 Tits 系统。
\item
  令 \(\mathbf{H} = \mathbf{SL}(2, \mathbf{A})\) 为 G 的子群，由系数属于 A 的矩阵组成。证明 \((\mathbf{H}, \mathbf{B} \cap \mathbf{H}, \mathbf{N}, \{s\})\) 是一个 Tits 系统。与习题 18 e) 比较。
\item
  令 \(\hat{\mathbf{A}}\) 为 \(\mathbf{A}\) 的完备化，令 \(\hat{\mathbf{G}},\hat{\mathbf{B}},\hat{\mathbf{N}},\hat{\mathbf{T}}\) 为如上定义但将 \(\mathbf{A}\) 替换为 \(\hat{\mathbf{A}}\) 得到的群。证明 \(\mathbf{G}\) 到 \(\hat{\mathbf{G}}\) 的嵌入定义了一个同构，将建筑 \(\mathbf{I}\)（与 \((\mathbf{G},\mathbf{B})\) 相关）映到建筑 \(\hat{\mathbf{I}}\)（与 \((\hat{\mathbf{G}},\hat{\mathbf{B}})\) 相关）（习题 10）。令 \((\mathbf{I},\mathfrak{U})\)（分别为 \((\hat{\mathbf{I}},\hat{\mathfrak{U}})\)）为与 \((\mathbf{G},\mathbf{B},\mathbf{N})\)（分别为 \((\hat{\mathbf{G}},\hat{\mathbf{B}},\hat{\mathbf{N}})\)）相关的结构建筑（习题 12）：证明 \(j(\mathfrak{U})\subset \hat{\mathfrak{U}}\)，但 \(j(\mathfrak{U})\neq \hat{\mathfrak{U}}\)（若 \(\mathbf{A}\neq \hat{\mathbf{A}}\)）（注意，公寓 \(\hat{\mathfrak{U}}\)（分别为 \(\mathfrak{U}\)）与 \(\mathbf{T}\)（分别为 \(\hat{\mathbf{T}}\)）在 \(\mathbf{G}\)（分别为 \(\hat{\mathbf{G}}\)）下的共轭子群一一对应）。
\end{enumerate}

\begin{enumerate}
\def\labelenumi{\arabic{enumi})}
\setcounter{enumi}{21}
\tightlist
\item
  令 \(G\) 为群，\(B\) 为 \(G\) 的子群。
\end{enumerate}

\begin{enumerate}
\def\labelenumi{\alph{enumi})}
\tightlist
\item
  证明下列条件等价：
\end{enumerate}

\begin{enumerate}
\def\labelenumi{(\roman{enumi})}
\item
  \(B \cap gBg^{-1}\) 在 \(B\) 中对所有 \(g \in G\) 具有有限指数；
\item
  关于 B 的每个双陪集 BgB 都是关于 B 的左陪集的有限并。
\end{enumerate}

更精确地，证明对所有 \(g \in G\)，指数 \(q_{g}\)（\(B \cap gBg^{-1}\) 在 B 中的指数）等于双陪集 BgB 中所包含的关于 B 的左陪集数。证明 \(q_{gh} \leqslant q_{g}q_{h}\) 对所有 \(g, h \in G\) 成立。

从现在起假设条件 (i) 和 (ii) 均已满足，并记 k 为可交换环。对 \(t \in G/B\)（分别对 \(t \in B \backslash G/B\)），记 \(a_{t}\) 为从 G 到 k 的映射，其定义为 \(a_{t}(g) = 0\)（若 \(g \notin t\)）且 \(a_{t}(g) = 1\)（若 \(g \in t\)）。令 L（分别为 H）为由 \(a_{t}\)（其中 \(t \in G/B\)，分别 \(t \in B \backslash G/B\)）生成的 k-模。

\begin{enumerate}
\def\labelenumi{\alph{enumi})}
\setcounter{enumi}{1}
\item
  证明存在唯一的线性形式 \(\mu\)，定义在 \(\mathbf{L}\) 上，使得 \(\mu(a_t) = 1\) 对所有 \(t \in \mathbf{G} / \mathbf{B}\) 成立。
\item
  令 \(\varphi \in \mathbf{L}\) 且 \(\psi \in \mathbf{H}\)。证明对所有 \(x \in \mathbf{G}\)，映射
\end{enumerate}

\[
\theta_ {x}: y \mapsto \varphi (y) \psi (y ^ {- 1} x)
\]

属于 L，且映射 \(\varphi * \psi : x \mapsto \mu(\theta_x)\) 属于 L。此外，若 \(\varphi \in H\)，则 \(\varphi * \psi \in H\)。证明映射 \((\varphi, \psi) \mapsto \varphi * \psi\) 使 H 成为 \(k\) 上的代数，单位元为 \(a_B\)，并使 L 成为右 H-模。

代数 \(\mathbf{H}\) 称为 \(\mathbf{G}\) 关于 \(\mathbf{H}\) 的 Hecke 代数，记作 \(\mathrm{H}_k(\mathbf{G},\mathbf{B})\)。

\begin{enumerate}
\def\labelenumi{\alph{enumi})}
\setcounter{enumi}{3}
\tightlist
\item
  证明对 \(t, t' \in B \backslash G / B\)，
\end{enumerate}

\[
a _ {t} * a _ {t ^ {\prime}} = \sum_ {t ^ {\prime \prime}} m (t, t ^ {\prime}; t ^ {\prime \prime}) a _ {t ^ {\prime \prime}},
\]

其中，\(m(t, t'; t'')\) 是包含于 \(t \cap gt'^{-1}\) 的关于 B 的陪集数（对所有 \(g \in t''\) 均如此）。

\begin{enumerate}
\def\labelenumi{\alph{enumi})}
\setcounter{enumi}{4}
\tightlist
\item
  令 G 通过左平移作用于 L。证明 H 在 L 上的作用定义了从 H 到由此得到的 G 在 L 上的线性表示的交换子的同构。
\end{enumerate}

\(^{*}f\) ) 假设 G 有限，且 k 的特征不整除 G 的阶。证明 \(H_{k}(G,B)\) 在 k 上绝对半单（《代数》第 VIII 章，§7，第 5 项）（使用 Maschke 定理（第 V 章，附录）以及《代数》第 VIII 章 §5 第 1 项的命题 3）。

\begin{enumerate}
\def\labelenumi{\alph{enumi})}
\setcounter{enumi}{6}
\tightlist
\item
  假设 G 是拓扑群且 B 是 G 的紧开子群。证明条件 (i) 和 (ii) 得到满足，并且当 k = R 或 C 时，乘积 \(\varphi * \psi\) 正是相对于 G 上右 Haar 测度的卷积乘积，其归一化条件为 \(\mu(B) = 1\)（参见《积分》第 VIII 章，§4，第 5 项）。
\end{enumerate}
% semantic-fragments: legacy-input-seam
\relax{}
\begin{enumerate}
\def\labelenumi{\arabic{enumi})}
\setcounter{enumi}{22}
\tightlist
\item
  设 \((\mathrm{W},\mathrm{S})\) 是一个 Coxeter 系统，\(k\) 是交换环。假设对于所有 \(s\in S\) ，给定 \(\lambda_{s}\) 中的两个元素 \(\mu_s\) 和 \(k\)，使得当 \(\lambda_{s} = \lambda_{s'}\) 与 \(\mu_{s} = \mu_{s'}\) 在 W 中共轭时，\(s\) 且 \(s'\)。置 \(\mathbf{E} = k^{(\mathbf{W})}\)，并令 \(\{e_w\}\) 为 E 的标准基。
\end{enumerate}

\begin{enumerate}
\def\labelenumi{\alph{enumi})}
\tightlist
\item
  证明 \(E\) 有唯一的代数结构，使得对于 \(s \in S\) 和 \(w \in W\)，
\end{enumerate}

\[
e _ {s}. e _ {w} = \left\{ \begin{array}{l l} e _ {s w} & \text {if} l (s w) > l (w) \\ \lambda_ {s} e _ {w} + \mu_ {s} e _ {s w} & \text {if} l (s w) <   l (w) \end{array} \right.
\]

(引入 E 的端同态 \(P_{s}\)，其由上述公式定义，其中以 \(e_{s}.e_{w}\) 替换 \(\mathrm{P}_{s}(w)\)；再引入端同态 \(Q_{s}=jP_{s}j^{-1}\)，其中 j 表示由 \(j(e_{w})=e_{w^{-1}}\) 定义的 E 的自同态；注意到条件 \(P_{s}Q_{t}=Q_{t}P_{s}\) 和 \(s,t\in S\) 蕴含 sw=wt，从而证明对于 \(l(swt)=l(w)\) 有 \(l(sw)=l(wt)\)；然后仿照习题 3 c) 论证。赋予此代数结构的模 E 记为 \(\mathrm{E}_{k}((\lambda_{s}),(\mu_{s}))\)。证明 \(\mathrm{E}((0),(1))\) 是 W 的群代数 k{[}W{]}（《代数》，第 III 章，§2，第 6 号）。

b) 证明生成元族 \((e_s)_{s\in S}\) 以及关系

\[
e _ {s} ^ {2} = \lambda_ {s} e _ {s} + \mu_ {s} \mathrm{for} s \in S
\]

\((e_s e_t)^r = (e_t e_s)^r\) 其中 \(s, t \in S\) 且 \(st\) 的有限偶阶为 \(2r\) \((e_s e_t)^r e_s = (e_t e_s)^r e_t\) 其中 \(s, t \in S\) 且 \(st\) 的有限奇阶为 \(2r + 1\)

给出代数 \(\mathbf{E}\) 的一个表示（仿照§1 第 6 号定理 1 的证明）。

\begin{enumerate}
\def\labelenumi{\arabic{enumi})}
\setcounter{enumi}{23}
\tightlist
\item
  设 G 是群，B 是 G 的一个 Tits 子群（习题 3，以下沿用其中的记号）。假设对于所有 \(s \in S\)，双陪集 \(C(s)\) 是关于 B 的有限个左陪集（其数目为 \(q_{s}\)）的并。沿用习题 22 的记号，并对所有 \(a_{w} = a_{\mathrm{C}(w)}\) 置 \(w \in W\)。
\end{enumerate}

\begin{enumerate}
\def\labelenumi{\alph{enumi})}
\item
  证明习题 22 的条件 (i) 和 (ii) 得到满足。因此可以谈论 Hecke 代数 \(\mathrm{H}_k(\mathbf{G},\mathbf{B})\)（\(k\) 为交换环），其中 \((a_w)_{w\in \mathbb{W}}\) 是一组基。
\item
  设 \(s \in S\) 且 \(w \in W\)。证明 \(a_s * a_s = (q_s - 1)a_s + q_s\)；证明若 \(l(sw) > l(w)\)，则 \(a_s * a_w = a_{sw}\)。
\item
  证明从与 Coxeter 系统 (W,S) 相关的代数 \(\mathrm{E}_{k}((q_{s}-1),(q_{s}))\)（习题 23）到 \(\mathrm{H}_{k}(\mathrm{G},\mathrm{B})\) 的线性映射是代数同构，其中对所有 \(e_{w}\) 将 \(a_{w}\) 送到 \(w \in W\)。
\end{enumerate}

\begin{enumerate}
\def\labelenumi{\arabic{enumi})}
\setcounter{enumi}{24}
\tightlist
\item
  沿用习题 8 的记号。另假设对于所有 \(s \in S\)，指数 \(q_s\) 是 \(B \cap gBg^{-1}\) 在 \(B\) 中的指数，并且对所有 \(g \in BsB\) 该指数是有限的。
\end{enumerate}

\begin{enumerate}
\def\labelenumi{\alph{enumi})}
\item
  证明对 \((\tilde{\mathbf{G}},\mathbf{B})\) 满足习题 22 的条件 (i) 和 (ii)。
\item
  证明对于所有 \(\gamma \in \Gamma\)，映射 \(x \mapsto \gamma x \gamma^{-1}\) 定义 Hecke 代数 \(\sigma\) 的一个自同态 \(\mathrm{H}_k(\mathrm{G}, \mathrm{B})\)，且 \(\sigma\) 只依赖于 \(\omega\) 在 \(\gamma\) 中的类 \(\Omega = \Gamma / (\Gamma \cap \mathrm{B})\)。
\item
  令 \(k[\Omega]\) 为 \(\Omega\) 的群代数，\((e_{\omega})\) 为其标准基。证明从 \(j\) 到 \(k[\Omega] \otimes_k \mathrm{H}_k(\mathbf{G}, \mathbf{B})\) 的线性映射 \(\mathrm{H}_k(\tilde{\mathrm{G}}, \mathrm{B})\) 由 \(j(e_{\omega} \otimes a_{\mathrm{B}w\mathrm{B}}) = a_{\mathrm{B}\omega w\mathrm{B}}\) 定义（采用习题 22 的记号），对于 \(\omega \in \Omega\) 和 \(w \in W\)，是双射，并且
\end{enumerate}

\[
j ^ {- 1} (j (e _ {\omega} \otimes x) j (e _ {w ^ {\prime}} \otimes y)) = e _ {w w ^ {\prime}} \otimes \sigma_ {\omega} (x) y
\]

其中 \(\omega, \omega' \in \Omega\)，且 \(x, y \in \mathrm{H}_k(\mathrm{G}, \mathrm{B})\)。

\begin{enumerate}
\def\labelenumi{\arabic{enumi})}
\setcounter{enumi}{25}
\tightlist
\item
  若 E 是交换域 k 上有限秩的绝对半单代数，则 E 的数值不变量是整数序列 \((n_{1},\ldots,n_{r})\)，满足 \(n_{1}\geqslant\cdots\geqslant n_{r}>0\)，并且对于 k 的任意代数闭包 \(\bar{k}\otimes_{k}E\)，\(\bar{k}\) 同构于 \(\prod M_{n_{i}}(\bar{k})\)。
\end{enumerate}

设 V 是整环，K 是其分式域，\(\varphi\) 是从 V 到交换域 k 的同态，E 是 V-代数。假设 E 是有限秩自由 V-模，并置 \(E_{0} = E \otimes_{V} k\) 和 \(E_{1} = E \otimes_{V} K\)。

\begin{enumerate}
\def\labelenumi{\alph{enumi})}
\item
  假设 \((x,y)\mapsto\mathrm{Tr}_{\mathbb{E}_{0}/k}(xy)\) 上的双线性型 \(E_{0}\) 非退化。证明 \(E_{0}\) 和 \(E_{1}\) 都是绝对半单的（参见《代数》，第 IX 章，§2，习题 1）。
\item
  假设 \(E_{0}\) 和 \(E_{1}\) 分别是 k 和 K 上的绝对半单代数。证明 \(E_{0}\) 和 \(E_{1}\) 具有相同的数值不变量（化归到 k 和 K 都代数闭的情形）；令 \((e_{i})\) 为 E 在 V 上的一组基，令 \((\mathbf{X}_i)\) 为不定元。证明 \(\sum_{i}\mathbf{X}_{i}e_{i}\) 中元素 \(\mathrm{E}_1\otimes_{\mathrm{K}}\mathrm{K}[(\mathrm{X}_i)]\) 的特征多项式（分别地，\(\mathrm{E}_0\otimes_kk[(\mathrm{X}_i)]\) 中相应元素的特征多项式）具有形式 \(\mathrm{P} = \prod_j\mathrm{P}_j^{n_j}\)（分别为 \(\mathrm{Q} = \prod_k\mathrm{Q}_k^{m_k}\)），其中 \((n_1,\ldots ,n_r)\)（分别为 \((m_1,\dots ,m_s))\) 是 \(\mathrm{E}_1\)（分别为 \(\mathrm{E}_0\)）的数值不变量，且 \(\deg (\mathrm{P}_j) = n_j\)（分别为 \(\deg (\mathrm{Q}_k) = m_k\)）。证明 \(\mathrm{P}_j\in \mathrm{V}[(\mathrm{X}_i)]\) 且 \(\mathrm{Q} = \varphi (\mathrm{P})\)。推知存在整数 \(c_{jk}\geqslant 0\)，使得
\end{enumerate}

\[
m _ {k} = \sum_ {j} c _ {j k} n _ {j}, \text { and } n _ {j} = \sum_ {k} c _ {j k} m _ {k}.
\]

*27) 设 \((\mathrm{G},\mathrm{B},\mathrm{N},\mathrm{S})\) 是一个 Tits 系统，且 \(k\) 是交换域。假设 G 是有限群，且 \(k\) 的特征既不整除 G 的阶，也不整除 Weyl 群 \(\mathrm{W} = \mathrm{N} / (\mathrm{B}\cap \mathrm{N})\) 的阶。证明代数 \(\mathrm{H}_k(\mathrm{G},\mathrm{B})\)（习题 22）和 \(k[\mathrm{W}]\) 都是绝对半单的并具有相同的数值不变量，因而当 \(k\) 代数闭时二者同构（令 \(q_{s}\) 为 \(\mathrm{B}\cap g\mathrm{Bg}^{-1}\)、\(g\in \mathrm{BsB}\) 时 \(s\in S\) 在 B 中的指数）。考虑按照习题 23 从 Coxeter 系统 (W,S) 和多项式环 \(\mathrm{E}_{k[\mathrm{X}]}((\mathrm{X}(q_s - 1)),(1 + \mathrm{X}(q_s - 1)))\) 构造的代数 \(k[\mathrm{X}]\)，并使用习题 26 a) 和 b)，注意根据 Maschke 定理（第 V 章，附录），双线性型 \(\mathrm{Tr}_{k[\mathrm{W}] / k}(xy)\) 是非退化的）。*

\begin{enumerate}
\def\labelenumi{\arabic{enumi})}
\setcounter{enumi}{27}
\item
  设 G 是群，M 是 G 的极大子群，U 是 M 的正规子群，满足第 7 号的条件 (R)。假设 G 等于其换位子群，由 U 的共轭子群之并生成，且 M 的共轭子群之交化为单位元。证明 G 是单群。（考虑 G 的一个不同于 \{1\} 的正规子群 N；证明 G = N.M，进而证明 G = N.U。）
\item
  设 H 是非阿贝尔单群。设 \(\theta\) 是 H 的一个阶为素数 p 的自同构，U 是由 \(\theta\) 对应的 Z/pZ 与 H 的半直积。证明若 \(\theta\) 不是内自同构，则 U 的正规子群只有 \(\{1\}\)、H 和 U。推知 U 不可解，但满足第 7 号的条件 (R)。应用此结论证明对称群 \(S_{n}\) 满足条件 (R)。
\end{enumerate}

